% META: {"id": "1NSI-PM-EVAL-B", "chapitre": "1NSI-PROJET-METHODES", "type_objet": "evaluation", "version": "B", "capacites": ["P-HIST-01", "BO-PREAMBULE-DEMARCHE-DE-PROJET", "BO-PREAMBULE-COMPETENCES-METHODE", "BO-PREAMBULE-COMPETENCES-ORALES"], "mode_creation": "ex_nihilo", "sources_inspiration": [], "duree_min": 45, "status": "verified"}
\section*{Évaluation B --- Projet et méthodes}
\textit{Durée : 45 minutes. Ordinateur avec interpréteur Python autorisé.}

\baremeIndicatif{Ex. 1 : 8 pts (histoire, démarche de projet) ; Ex. 2 : 8 pts (débogage) ; Ex. 3 : 4 pts (documentation et oral)}

\bigskip

\textbf{Exercice 1} --- \textit{Histoire et gestion de projet} \hfill (8 points)

\begin{enumerate}
  \item (3 pts) En utilisant \lstinline{evenements_entre} du cours sur la chronologie \lstinline{[(1936, ...), (1969, ...), (1991, ...)]}, lister les événements situés entre 1960 et 1991.
  \item (2 pts) Un projet est découpé en $5$ jalons, dont $3$ sont terminés. Calculer son avancement avec \lstinline{avancement} du cours.
  \item (3 pts) Le programme officiel précise que les professeurs veillent à ce que les projets restent « d'une ambition raisonnable ». Pourquoi cette recommandation est-elle importante pour une équipe de $2$ à $4$ élèves ?
\end{enumerate}

% BEGIN-VERIFY
% evenements = [
%     (1936, "Concept de machine universelle", "Alan Turing"),
%     (1969, "Debuts du reseau Internet", "DARPA"),
%     (1991, "Premiere version publique de Python", "Guido van Rossum"),
% ]
% def evenements_entre(evenements, debut, fin):
%     return [e for e in evenements if debut <= e[0] <= fin]
% resultat = evenements_entre(evenements, 1960, 1991)
% assert [e[0] for e in resultat] == [1969, 1991]
% jalons = [
%     {"nom": "Maquette", "termine": True},
%     {"nom": "Backend", "termine": True},
%     {"nom": "Frontend", "termine": True},
%     {"nom": "Documentation", "termine": False},
%     {"nom": "Presentation", "termine": False},
% ]
% def avancement(jalons):
%     return sum(1 for j in jalons if j["termine"]) / len(jalons) * 100
% assert avancement(jalons) == 60.0
% END-VERIFY

\bigskip

\textbf{Exercice 2} --- \textit{Débogage} \hfill (8 points)

\begin{python}
def acceder_element_bugue(liste, indice):
    return liste[indice]

acceder_element_bugue([1, 2, 3], 5)
\end{python}

\begin{enumerate}
  \item (3 pts) Quel type d'erreur ce programme provoque-t-il ? Quel est le message exact ?
  \item (3 pts) Écrire une version corrigée \lstinline{acceder_element(liste, indice)} qui utilise une précondition (\lstinline{assert}) pour vérifier que l'indice est valide, avec un message clair.
  \item (2 pts) Pourquoi une précondition explicite est-elle préférable à laisser l'erreur native (\lstinline{IndexError}) se produire sans explication ?
\end{enumerate}

% BEGIN-VERIFY
% def acceder_element_bugue(liste, indice):
%     return liste[indice]
% erreur_native = False
% try:
%     acceder_element_bugue([1, 2, 3], 5)
% except IndexError as e:
%     erreur_native = True
%     message = str(e)
% assert erreur_native is True
% assert message == "list index out of range"
% def acceder_element(liste, indice):
%     assert 0 <= indice < len(liste), "indice hors des bornes de la liste"
%     return liste[indice]
% assert acceder_element([1, 2, 3], 1) == 2
% erreur_precondition = False
% try:
%     acceder_element([1, 2, 3], 5)
% except AssertionError:
%     erreur_precondition = True
% assert erreur_precondition is True
% END-VERIFY

\bigskip

\textbf{Exercice 3} --- \textit{Documentation et presentation orale} \hfill (4 points)

On donne cette fonction, ecrite pour un projet de gestion de notes :

\begin{python}
def moyenne(notes, coefficients):
    total = 0
    poids = 0
    for note, coefficient in zip(notes, coefficients):
        total += note * coefficient
        poids += coefficient
    return total / poids
\end{python}

\begin{enumerate}
  \item (2 pts) Ajouter a cette fonction une docstring conforme aux usages du cours : role,
  paramètres avec leur type, valeur renvoyee, et la precondition qui evite la division par
  zero. On ecrira la fonction complete.
  \item (2 pts) Vous presentez ce projet a l'oral en trois minutes. Rediger le plan de votre
  intervention en trois points, et indiquer pour chacun l'argument principal que vous
  avancerez. Un des trois points devra justifier un choix technique et non le decrire.
\end{enumerate}

% BEGIN-VERIFY
% def moyenne(notes, coefficients):
%     """Renvoie la moyenne ponderee des notes.
%
%     notes : liste de nombres, les notes obtenues.
%     coefficients : liste de nombres positifs, de meme longueur que notes.
%     Renvoie un nombre : la moyenne ponderee.
%     Precondition : la somme des coefficients est strictement positive.
%     """
%     assert sum(coefficients) > 0, "la somme des coefficients doit etre non nulle"
%     total = 0
%     poids = 0
%     for note, coefficient in zip(notes, coefficients):
%         total += note * coefficient
%         poids += coefficient
%     return total / poids
%
% assert moyenne([10, 20], [1, 1]) == 15
% assert moyenne([12, 8], [3, 1]) == 11
% # La docstring est bien presente et nomme les quatre elements attendus.
% doc = moyenne.__doc__
% for attendu in ("notes", "coefficients", "Renvoie", "Precondition"):
%     assert attendu in doc
% # La precondition transforme une division par zero en message explicite.
% try:
%     moyenne([10], [0])
% except AssertionError as e:
%     assert "coefficients" in str(e)
% else:
%     raise AssertionError("la precondition n'a pas ete declenchee")
% END-VERIFY
